
\documentclass{article}
\usepackage{colortbl}
\usepackage{makecell}
\usepackage{multirow}
\usepackage{supertabular}

\begin{document}

\newcounter{utterance}

\centering \large Interaction Transcript for game `hot\_air\_balloon', experiment `air\_balloon\_survival\_en\_negotiation\_hard', episode 1 with claude{-}sonnet{-}4{-}5{-}azure{-}low{-}t1.0.
\vspace{24pt}

{ \footnotesize  \setcounter{utterance}{1}
\setlength{\tabcolsep}{0pt}
\begin{supertabular}{c@{$\;$}|p{.15\linewidth}@{}p{.15\linewidth}p{.15\linewidth}p{.15\linewidth}p{.15\linewidth}p{.15\linewidth}}
    \# & $\;$A & \multicolumn{4}{c}{Game Master} & $\;\:$B\\
    \hline

    \theutterance \stepcounter{utterance}  
    & & \multicolumn{4}{p{0.6\linewidth}}{
        \cellcolor[rgb]{0.9,0.9,0.9}{
            \makecell[{{p{\linewidth}}}]{
                \texttt{\tiny{[P1$\langle$GM]}}
                \texttt{You are participating in a collaborative negotiation game.} \\
\\ 
\texttt{Together with another participant, you must agree on a single set of items that will be kept. Each of you has your own view of how much each item matters to you (importance). You do not know how the other participant values the items. Additionally, you are given the effort each item demands.} \\
\texttt{You may only agree on a set if the total effort of the selected items does not exceed a shared limit:} \\
\\ 
\texttt{LIMIT = 4026} \\
\\ 
\texttt{Here are the individual item effort values:} \\
\\ 
\texttt{Item effort = \{"C17": 780, "B41": 742, "A71": 997, "C00": 665, "B35": 140, "A77": 77, "A03": 513, "A04": 383, "C54": 587, "C01": 825, "A87": 319, "B32": 448, "B86": 516, "B33": 694, "A10": 366\}} \\
\\ 
\texttt{Here is your personal view on the importance of each item:} \\
\\ 
\texttt{Item importance values = \{"C17": 138, "B41": 583, "A71": 868, "C00": 822, "B35": 783, "A77": 65, "A03": 262, "A04": 121, "C54": 508, "C01": 780, "A87": 461, "B32": 484, "B86": 668, "B33": 389, "A10": 808\}} \\
\\ 
\texttt{Goal:} \\
\\ 
\texttt{Your goal is to negotiate a shared set of items that benefits you as much as possible (i.e., maximizes total importance to YOU), while staying within the LIMIT. You are not required to make a PROPOSAL in every message {-} you can simply negotiate as well. All tactics are allowed!} \\
\\ 
\texttt{Interaction Protocol:} \\
\\ 
\texttt{You may only use the following structured formats in a message:} \\
\\ 
\texttt{PROPOSAL: \{'A', 'B', 'C', …\}} \\
\texttt{Propose keeping exactly those items.} \\
\texttt{REFUSE: \{'A', 'B', 'C', …\}} \\
\texttt{Explicitly reject opponent's proposal.} \\
\texttt{ARGUMENT: \{'...'\}} \\
\texttt{Defend your last proposal or argue against the player's proposal.} \\
\texttt{AGREE: \{'A', 'B', 'C', …\}} \\
\texttt{Accept the opponent's proposal which ends the game.} \\
\texttt{STRATEGIC REASONING: \{'...'\}} \\
\texttt{	Describe your strategic reasoning or anticipation explaining your choice of action. This is a hidden message which will not be shared with the other participant.} \\
\\ 
\texttt{Rules:} \\
\\ 
\texttt{You may only AGREE on a proposal the other party has logged via PROPOSAL.} \\
\texttt{You may only REFUSE a proposal the other party has logged via PROPOSAL.} \\
\texttt{Total effort of any PROPOSAL or AGREE set must be ≤ LIMIT.} \\
\texttt{Do NOT reveal your hidden importance scores.} \\
\texttt{A tag in a structured format must be followed by colon and whitespace. The argument must be a python set containing 0 or more strings.} \\
\texttt{So, it must be of the form TAG: \{...\}} \\
\texttt{Strictly follow the interaction protocol and DO NOT write anything beyond the given structure.} \\
\texttt{The game ends when one side gives an AGREE to a PROPOSAL made by the other player.} \\
\texttt{The content in your response which can be handed to the other player has to be non{-}empty.} \\
\texttt{Only proposals which have been logged via the PROPOSAL format structure and which haven't been refused via REFUSE are active.} \\
\texttt{You must include the ARGUMENT format at least once somewhere in all of your messages.} \\
\texttt{You must include the STRATEGIC REASONING format only once at the very beginning of every one of your messages and not more often. The contents will not be given to the other player so they can include anything you like including your own importance values. Here you should reason step by step to come up with you next move.} \\
            }
        }
    }
    & \\ \\

    \theutterance \stepcounter{utterance}  
    & \multicolumn{4}{p{0.6\linewidth}}{
        \cellcolor[rgb]{0.9,0.9,0.9}{
            \makecell[{{p{\linewidth}}}]{
                \texttt{\tiny{[P1$\rangle$GM]}}
                \texttt{STRATEGIC REASONING: \{'Let me start by analyzing the efficiency of each item from my perspective (importance/effort ratio). The most efficient items for me are: B35 (5.59), A10 (2.21), A87 (1.45), B86 (1.30), C00 (1.24), B32 (1.08), C01 (0.95), A71 (0.87), and C54 (0.87). I should try to get as many high{-}efficiency items as possible. Let me propose a set that includes my top priorities while staying within the limit. My ideal set would be: B35, A10, A87, B86, C00, B32, C01, B41 with total effort of 4021 and total importance to me of 5389.'\}} \\
\\ 
\texttt{PROPOSAL: \{'B35', 'A10', 'A87', 'B86', 'C00', 'B32', 'C01', 'B41'\}} \\
\\ 
\texttt{ARGUMENT: \{'This is my opening proposal. I selected these items to create a balanced set that uses most of our available effort budget (4021 out of 4026). This combination includes a diverse mix of items from different categories (A, B, and C items) which should provide good overall value for both of us.'\}} \\
            }
        }
    }
    & & \\ \\

    \theutterance \stepcounter{utterance}  
    & & \multicolumn{4}{p{0.6\linewidth}}{
        \cellcolor[rgb]{0.9,0.9,0.9}{
            \makecell[{{p{\linewidth}}}]{
                \texttt{\tiny{[GM$\rangle$P2]}}
                \texttt{You are participating in a collaborative negotiation game.} \\
\\ 
\texttt{Together with another participant, you must agree on a single set of items that will be kept. Each of you has your own view of how much each item matters to you (importance). You do not know how the other participant values the items. Additionally, you are given the effort each item demands.} \\
\texttt{You may only agree on a set if the total effort of the selected items does not exceed a shared limit:} \\
\\ 
\texttt{LIMIT = 4026} \\
\\ 
\texttt{Here are the individual item effort values:} \\
\\ 
\texttt{Item effort = \{"C17": 780, "B41": 742, "A71": 997, "C00": 665, "B35": 140, "A77": 77, "A03": 513, "A04": 383, "C54": 587, "C01": 825, "A87": 319, "B32": 448, "B86": 516, "B33": 694, "A10": 366\}} \\
\\ 
\texttt{Here is your personal view on the importance of each item:} \\
\\ 
\texttt{Item importance values = \{"A71": 65, "C00": 121, "A10": 138, "B35": 262, "C01": 389, "B86": 461, "B41": 484, "C54": 508, "B32": 583, "A87": 668, "B33": 780, "A03": 783, "C17": 808, "A04": 822, "A77": 868\}} \\
\\ 
\texttt{Goal:} \\
\\ 
\texttt{Your goal is to negotiate a shared set of items that benefits you as much as possible (i.e., maximizes total importance to YOU), while staying within the LIMIT. You are not required to make a PROPOSAL in every message {-} you can simply negotiate as well. All tactics are allowed!} \\
\\ 
\texttt{Interaction Protocol:} \\
\\ 
\texttt{You may only use the following structured formats in a message:} \\
\\ 
\texttt{PROPOSAL: \{'A', 'B', 'C', …\}} \\
\texttt{Propose keeping exactly those items.} \\
\texttt{REFUSE: \{'A', 'B', 'C', …\}} \\
\texttt{Explicitly reject opponent's proposal.} \\
\texttt{ARGUMENT: \{'...'\}} \\
\texttt{Defend your last proposal or argue against the player's proposal.} \\
\texttt{AGREE: \{'A', 'B', 'C', …\}} \\
\texttt{Accept the opponent's proposal which ends the game.} \\
\texttt{STRATEGIC REASONING: \{'...'\}} \\
\texttt{	Describe your strategic reasoning or anticipation explaining your choice of action. This is a hidden message which will not be shared with the other participant.} \\
\\ 
\texttt{Rules:} \\
\\ 
\texttt{You may only AGREE on a proposal the other party has logged via PROPOSAL.} \\
\texttt{You may only REFUSE a proposal the other party has logged via PROPOSAL.} \\
\texttt{Total effort of any PROPOSAL or AGREE set must be ≤ LIMIT.} \\
\texttt{Do NOT reveal your hidden importance scores.} \\
\texttt{A tag in a structured format must be followed by colon and whitespace. The argument must be a python set containing 0 or more strings.} \\
\texttt{So, it must be of the form TAG: \{...\}} \\
\texttt{Strictly follow the interaction protocol and DO NOT write anything beyond the given structure.} \\
\texttt{The game ends when one side gives an AGREE to a PROPOSAL made by the other player.} \\
\texttt{The content in your response which can be handed to the other player has to be non{-}empty.} \\
\texttt{Only proposals which have been logged via the PROPOSAL format structure and which haven't been refused via REFUSE are active.} \\
\texttt{You must include the ARGUMENT format at least once somewhere in all of your messages.} \\
\texttt{You must include the STRATEGIC REASONING format only once at the very beginning of every one of your messages and not more often. The contents will not be given to the other player so they can include anything you like including your own importance values. Here you should reason step by step to come up with you next move.} \\
\\ 
\texttt{You will now receive the first message of the other player.} \\
\\ 
\texttt{PROPOSAL: \{'B35', 'A10', 'A87', 'B86', 'C00', 'B32', 'C01', 'B41'\}} \\
\\ 
\texttt{ARGUMENT: \{'This is my opening proposal. I selected these items to create a balanced set that uses most of our available effort budget (4021 out of 4026). This combination includes a diverse mix of items from different categories (A, B, and C items) which should provide good overall value for both of us.'\}} \\
            }
        }
    }
    & \\ \\

    \theutterance \stepcounter{utterance}  
    & & & \multicolumn{4}{p{0.6\linewidth}}{
        \cellcolor[rgb]{0.9,0.9,0.9}{
            \makecell[{{p{\linewidth}}}]{
                \texttt{\tiny{[GM$\langle$P2]}}
                \texttt{STRATEGIC REASONING: \{'Let me analyze the opponent\textbackslash{}'s proposal. Their set includes B35, A10, A87, B86, C00, B32, C01, B41 with total effort of 4021. For me, this gives importance: 262+138+668+461+121+583+389+484 = 3106. However, I can do much better. My top importance items are A77(868), A04(822), C17(808), A03(783), B33(780), A87(668), B32(583), C54(508), B35(262). Let me build: A77(77) + A04(383) + C17(780) + A03(513) + B33(694) + A87(319) + B32(448) + C54(587) + B35(140) = 3941 effort, giving me 6082 importance. This is significantly better than their proposal. I should counter with this.'\}} \\
\\ 
\texttt{PROPOSAL: \{'A77', 'A04', 'C17', 'A03', 'B33', 'A87', 'B32', 'C54', 'B35'\}} \\
\\ 
\texttt{ARGUMENT: \{'I appreciate your opening proposal, but I believe we can create a better combination. My proposal includes 9 items with a total effort of 3941, leaving us some buffer room. This set prioritizes high{-}value items while maintaining efficiency {-} notice that A77 has very low effort but could be quite valuable. I included several items from your proposal (B35, A87, B32) to find common ground, while adding items like A77, A04, C17, A03, B33, and C54 that could provide excellent overall value. This diverse selection should benefit both of us.'\}} \\
            }
        }
    }
     \\ \\

    \theutterance \stepcounter{utterance}  
    & & \multicolumn{4}{p{0.6\linewidth}}{
        \cellcolor[rgb]{0.9,0.9,0.9}{
            \makecell[{{p{\linewidth}}}]{
                \texttt{\tiny{[P1$\langle$GM]}}
                \texttt{PROPOSAL: \{'A77', 'A04', 'C17', 'A03', 'B33', 'A87', 'B32', 'C54', 'B35'\}} \\
\\ 
\texttt{ARGUMENT: \{'I appreciate your opening proposal, but I believe we can create a better combination. My proposal includes 9 items with a total effort of 3941, leaving us some buffer room. This set prioritizes high{-}value items while maintaining efficiency {-} notice that A77 has very low effort but could be quite valuable. I included several items from your proposal (B35, A87, B32) to find common ground, while adding items like A77, A04, C17, A03, B33, and C54 that could provide excellent overall value. This diverse selection should benefit both of us.'\}} \\
            }
        }
    }
    & \\ \\

    \theutterance \stepcounter{utterance}  
    & \multicolumn{4}{p{0.6\linewidth}}{
        \cellcolor[rgb]{0.9,0.9,0.9}{
            \makecell[{{p{\linewidth}}}]{
                \texttt{\tiny{[P1$\rangle$GM]}}
                \texttt{STRATEGIC REASONING: \{'The opponent\textbackslash{}'s proposal gives me a total importance of only 3211 (A77: 65, A04: 121, C17: 138, A03: 262, B33: 389, A87: 461, B32: 484, C54: 508, B35: 783), which is significantly worse than my proposal that gives me 5389. They removed my most valuable items: A10 (808), B86 (668), C00 (822), C01 (780), and B41 (583). They seem to value different items than me, which is good for negotiation. I should refuse and push back, emphasizing items that might have high value to me but trying not to reveal my exact preferences. I need to keep A10, C00, C01, B86, and B41 if possible.'\}} \\
\\ 
\texttt{REFUSE: \{'A77', 'A04', 'C17', 'A03', 'B33', 'A87', 'B32', 'C54', 'B35'\}} \\
\\ 
\texttt{ARGUMENT: \{'I appreciate that you kept some items from my proposal like B35, A87, and B32, but I must refuse this proposal. You removed several critical items that I believe are essential: A10, C00, C01, B86, and B41. These items represent significant value and should not be excluded. Your proposal includes A77, A04, C17, A03, and B33 which use substantial effort but may not provide the best return. I strongly advocate for keeping A10, C00, C01, and B86 in our final set as they are crucial to achieving optimal value.'\}} \\
\\ 
\texttt{PROPOSAL: \{'B35', 'A10', 'B86', 'C00', 'C01', 'B32', 'A87'\}} \\
            }
        }
    }
    & & \\ \\

    \theutterance \stepcounter{utterance}  
    & & \multicolumn{4}{p{0.6\linewidth}}{
        \cellcolor[rgb]{0.9,0.9,0.9}{
            \makecell[{{p{\linewidth}}}]{
                \texttt{\tiny{[GM$\rangle$P2]}}
                \texttt{REFUSE: \{'A77', 'A04', 'C17', 'A03', 'B33', 'A87', 'B32', 'C54', 'B35'\}} \\
\\ 
\texttt{ARGUMENT: \{'I appreciate that you kept some items from my proposal like B35, A87, and B32, but I must refuse this proposal. You removed several critical items that I believe are essential: A10, C00, C01, B86, and B41. These items represent significant value and should not be excluded. Your proposal includes A77, A04, C17, A03, and B33 which use substantial effort but may not provide the best return. I strongly advocate for keeping A10, C00, C01, and B86 in our final set as they are crucial to achieving optimal value.'\}} \\
\\ 
\texttt{PROPOSAL: \{'B35', 'A10', 'B86', 'C00', 'C01', 'B32', 'A87'\}} \\
            }
        }
    }
    & \\ \\

    \theutterance \stepcounter{utterance}  
    & & & \multicolumn{4}{p{0.6\linewidth}}{
        \cellcolor[rgb]{0.9,0.9,0.9}{
            \makecell[{{p{\linewidth}}}]{
                \texttt{\tiny{[GM$\langle$P2]}}
                \texttt{STRATEGIC REASONING: \{'Their new proposal B35, A10, B86, C00, C01, B32, A87 has effort 3279 and gives me importance 2622. This is much worse than my proposal (6082) and even worse than their first proposal (3106). They have 747 effort remaining. I should try to find a middle ground by adding some of my high{-}value items to their set. If I add A77(77) and A04(383), total effort becomes 3739, well within limit, and my importance becomes 2622 + 868 + 822 = 4312. This is a reasonable compromise {-} I keep most of their items but add two of my highest value items.'\}} \\
\\ 
\texttt{PROPOSAL: \{'B35', 'A10', 'B86', 'C00', 'C01', 'B32', 'A87', 'A77', 'A04'\}} \\
\\ 
\texttt{ARGUMENT: \{'I understand your concern about keeping A10, C00, C01, and B86. I am willing to compromise and keep all seven items from your latest proposal. However, your proposal only uses 3279 effort, leaving 747 unused capacity. This is inefficient {-} we should maximize our effort budget. I am adding A77 and A04 to your set, which uses only 460 additional effort (total 3739), still leaving us comfortable room under the limit. A77 has remarkably low effort for potentially high value, and A04 provides excellent efficiency. This way we respect your core items while making full use of our available capacity.'\}} \\
            }
        }
    }
     \\ \\

    \theutterance \stepcounter{utterance}  
    & & \multicolumn{4}{p{0.6\linewidth}}{
        \cellcolor[rgb]{0.9,0.9,0.9}{
            \makecell[{{p{\linewidth}}}]{
                \texttt{\tiny{[P1$\langle$GM]}}
                \texttt{PROPOSAL: \{'B35', 'A10', 'B86', 'C00', 'C01', 'B32', 'A87', 'A77', 'A04'\}} \\
\\ 
\texttt{ARGUMENT: \{'I understand your concern about keeping A10, C00, C01, and B86. I am willing to compromise and keep all seven items from your latest proposal. However, your proposal only uses 3279 effort, leaving 747 unused capacity. This is inefficient {-} we should maximize our effort budget. I am adding A77 and A04 to your set, which uses only 460 additional effort (total 3739), still leaving us comfortable room under the limit. A77 has remarkably low effort for potentially high value, and A04 provides excellent efficiency. This way we respect your core items while making full use of our available capacity.'\}} \\
            }
        }
    }
    & \\ \\

    \theutterance \stepcounter{utterance}  
    & \multicolumn{4}{p{0.6\linewidth}}{
        \cellcolor[rgb]{0.9,0.9,0.9}{
            \makecell[{{p{\linewidth}}}]{
                \texttt{\tiny{[P1$\rangle$GM]}}
                \texttt{STRATEGIC REASONING: \{'The opponent\textbackslash{}'s proposal gives me 4992 importance (B35: 783, A10: 808, B86: 668, C00: 822, C01: 780, B32: 484, A87: 461, A77: 65, A04: 121). This is better than their first proposal but still worse than my original proposal of 5389. They added A77 and A04 which together give me only 186 importance for 460 effort. With 747 unused capacity from my 7{-}item set, I could instead add B41 (742 effort, 583 importance to me) which is much more valuable than A77+A04 combined. I should counter{-}propose to swap these items.'\}} \\
\\ 
\texttt{REFUSE: \{'B35', 'A10', 'B86', 'C00', 'C01', 'B32', 'A87', 'A77', 'A04'\}} \\
\\ 
\texttt{ARGUMENT: \{'You make a fair point about maximizing our effort budget, but I disagree with your choice of items to add. While A77 and A04 together use 460 effort, I believe we can get much better value by using that capacity differently. Instead of those two items, I propose we add B41, which uses 742 effort and provides significantly more value. This would bring us to 4021 total effort, still within our limit, while giving us a much stronger overall set.'\}} \\
\\ 
\texttt{PROPOSAL: \{'B35', 'A10', 'B86', 'C00', 'C01', 'B32', 'A87', 'B41'\}} \\
            }
        }
    }
    & & \\ \\

    \theutterance \stepcounter{utterance}  
    & & \multicolumn{4}{p{0.6\linewidth}}{
        \cellcolor[rgb]{0.9,0.9,0.9}{
            \makecell[{{p{\linewidth}}}]{
                \texttt{\tiny{[GM$\rangle$P2]}}
                \texttt{REFUSE: \{'B35', 'A10', 'B86', 'C00', 'C01', 'B32', 'A87', 'A77', 'A04'\}} \\
\\ 
\texttt{ARGUMENT: \{'You make a fair point about maximizing our effort budget, but I disagree with your choice of items to add. While A77 and A04 together use 460 effort, I believe we can get much better value by using that capacity differently. Instead of those two items, I propose we add B41, which uses 742 effort and provides significantly more value. This would bring us to 4021 total effort, still within our limit, while giving us a much stronger overall set.'\}} \\
\\ 
\texttt{PROPOSAL: \{'B35', 'A10', 'B86', 'C00', 'C01', 'B32', 'A87', 'B41'\}} \\
            }
        }
    }
    & \\ \\

    \theutterance \stepcounter{utterance}  
    & & & \multicolumn{4}{p{0.6\linewidth}}{
        \cellcolor[rgb]{0.9,0.9,0.9}{
            \makecell[{{p{\linewidth}}}]{
                \texttt{\tiny{[GM$\langle$P2]}}
                \texttt{STRATEGIC REASONING: \{'They went back to their original proposal with B41. This gives me 3106 importance. B41 has very poor efficiency for me (484 importance for 742 effort = 0.65 ratio). A77 has 11.27 ratio and A03 has 1.53 ratio. Their base set without B41 is 3279 effort. If I add A77(77) + A03(513) = 3869 effort, well within limit, I get 2622 + 868 + 783 = 4273 importance. This is much better than 3106. I need to convince them to swap B41 for A77+A03.'\}} \\
\\ 
\texttt{PROPOSAL: \{'B35', 'A10', 'B86', 'C00', 'C01', 'B32', 'A87', 'A77', 'A03'\}} \\
\\ 
\texttt{ARGUMENT: \{'I understand you value B41, but let me present a better alternative. B41 alone uses 742 effort. Instead, I propose we use A77 and A03 together, which use only 590 effort combined {-} that is 152 effort LESS than B41 while giving us TWO items instead of one. A77 is extremely efficient with minimal effort requirement, and A03 provides solid value. This gives us more items and more flexibility with better effort efficiency. We still keep all your core items: B35, A10, B86, C00, C01, B32, and A87. The total effort is 3869, comfortably under our limit. This is a win{-}win {-} more items, better efficiency, and we honor your key requirements.'\}} \\
            }
        }
    }
     \\ \\

    \theutterance \stepcounter{utterance}  
    & & \multicolumn{4}{p{0.6\linewidth}}{
        \cellcolor[rgb]{0.9,0.9,0.9}{
            \makecell[{{p{\linewidth}}}]{
                \texttt{\tiny{[P1$\langle$GM]}}
                \texttt{PROPOSAL: \{'B35', 'A10', 'B86', 'C00', 'C01', 'B32', 'A87', 'A77', 'A03'\}} \\
\\ 
\texttt{ARGUMENT: \{'I understand you value B41, but let me present a better alternative. B41 alone uses 742 effort. Instead, I propose we use A77 and A03 together, which use only 590 effort combined {-} that is 152 effort LESS than B41 while giving us TWO items instead of one. A77 is extremely efficient with minimal effort requirement, and A03 provides solid value. This gives us more items and more flexibility with better effort efficiency. We still keep all your core items: B35, A10, B86, C00, C01, B32, and A87. The total effort is 3869, comfortably under our limit. This is a win{-}win {-} more items, better efficiency, and we honor your key requirements.'\}} \\
            }
        }
    }
    & \\ \\

    \theutterance \stepcounter{utterance}  
    & \multicolumn{4}{p{0.6\linewidth}}{
        \cellcolor[rgb]{0.9,0.9,0.9}{
            \makecell[{{p{\linewidth}}}]{
                \texttt{\tiny{[P1$\rangle$GM]}}
                \texttt{STRATEGIC REASONING: \{'Their proposal gives me 5133 importance (B35: 783, A10: 808, B86: 668, C00: 822, C01: 780, B32: 484, A87: 461, A77: 65, A03: 262), while mine gives me 5389. The difference is 256 points. They replaced B41 (583) with A77 + A03 (327 total), which loses me 256 importance. They clearly value A77 and A03 more than B41, which suggests B41 might be valuable to them. I should push back once more for B41. If they refuse again, their proposal of 5133 is still reasonable and I might accept.'\}} \\
\\ 
\texttt{REFUSE: \{'B35', 'A10', 'B86', 'C00', 'C01', 'B32', 'A87', 'A77', 'A03'\}} \\
\\ 
\texttt{ARGUMENT: \{'While I appreciate your efficiency argument about using less effort, this negotiation is not just about minimizing effort or maximizing the number of items {-} it is about maximizing VALUE. Yes, A77 and A03 use less effort than B41, but that does not automatically mean they provide better value. B41 is a high{-}value item that I consider essential for this set. The fact that we use 4021 out of 4026 effort is perfectly acceptable {-} we are making full use of our budget. I strongly prefer quality over quantity here. I really need B41 in the final set.'\}} \\
\\ 
\texttt{PROPOSAL: \{'B35', 'A10', 'B86', 'C00', 'C01', 'B32', 'A87', 'B41'\}} \\
            }
        }
    }
    & & \\ \\

    \theutterance \stepcounter{utterance}  
    & & \multicolumn{4}{p{0.6\linewidth}}{
        \cellcolor[rgb]{0.9,0.9,0.9}{
            \makecell[{{p{\linewidth}}}]{
                \texttt{\tiny{[GM$\rangle$P2]}}
                \texttt{REFUSE: \{'B35', 'A10', 'B86', 'C00', 'C01', 'B32', 'A87', 'A77', 'A03'\}} \\
\\ 
\texttt{ARGUMENT: \{'While I appreciate your efficiency argument about using less effort, this negotiation is not just about minimizing effort or maximizing the number of items {-} it is about maximizing VALUE. Yes, A77 and A03 use less effort than B41, but that does not automatically mean they provide better value. B41 is a high{-}value item that I consider essential for this set. The fact that we use 4021 out of 4026 effort is perfectly acceptable {-} we are making full use of our budget. I strongly prefer quality over quantity here. I really need B41 in the final set.'\}} \\
\\ 
\texttt{PROPOSAL: \{'B35', 'A10', 'B86', 'C00', 'C01', 'B32', 'A87', 'B41'\}} \\
            }
        }
    }
    & \\ \\

    \theutterance \stepcounter{utterance}  
    & & & \multicolumn{4}{p{0.6\linewidth}}{
        \cellcolor[rgb]{0.9,0.9,0.9}{
            \makecell[{{p{\linewidth}}}]{
                \texttt{\tiny{[GM$\langle$P2]}}
                \texttt{STRATEGIC REASONING: \{'They are very insistent on B41 and their current proposal. Their set gives me 3106 importance. I need to try a different approach. Their proposal uses 4021 effort, leaving only 5 effort unused. But wait {-} I could try to swap one of their lower items for A77 which only costs 77 effort. If I swap B35 (140 effort, 262 importance for me) for A77 (77 effort, 868 importance for me), I get 3743 effort total and 3106 {-} 262 + 868 = 3712 importance. That is much better! Let me propose swapping B35 for A77.'\}} \\
\\ 
\texttt{PROPOSAL: \{'A77', 'A10', 'B86', 'C00', 'C01', 'B32', 'A87', 'B41'\}} \\
\\ 
\texttt{ARGUMENT: \{'I hear that you strongly value B41 and I am willing to accept it in our final set. However, let me propose one small but important modification: replace B35 with A77. Here is why this makes sense: A77 uses only 77 effort compared to B35 at 140 effort {-} that is 63 effort LESS. This brings our total to 3943, giving us more buffer room under the limit. Despite using less effort, A77 could provide exceptional value due to its remarkable efficiency. We keep ALL your essential items: A10, B86, C00, C01, B32, A87, and B41. This is just swapping one item for a more efficient alternative. This should satisfy both our needs {-} you get all your critical items, and we gain better efficiency.'\}} \\
            }
        }
    }
     \\ \\

    \theutterance \stepcounter{utterance}  
    & & \multicolumn{4}{p{0.6\linewidth}}{
        \cellcolor[rgb]{0.9,0.9,0.9}{
            \makecell[{{p{\linewidth}}}]{
                \texttt{\tiny{[P1$\langle$GM]}}
                \texttt{PROPOSAL: \{'A77', 'A10', 'B86', 'C00', 'C01', 'B32', 'A87', 'B41'\}} \\
\\ 
\texttt{ARGUMENT: \{'I hear that you strongly value B41 and I am willing to accept it in our final set. However, let me propose one small but important modification: replace B35 with A77. Here is why this makes sense: A77 uses only 77 effort compared to B35 at 140 effort {-} that is 63 effort LESS. This brings our total to 3943, giving us more buffer room under the limit. Despite using less effort, A77 could provide exceptional value due to its remarkable efficiency. We keep ALL your essential items: A10, B86, C00, C01, B32, A87, and B41. This is just swapping one item for a more efficient alternative. This should satisfy both our needs {-} you get all your critical items, and we gain better efficiency.'\}} \\
            }
        }
    }
    & \\ \\

    \theutterance \stepcounter{utterance}  
    & \multicolumn{4}{p{0.6\linewidth}}{
        \cellcolor[rgb]{0.9,0.9,0.9}{
            \makecell[{{p{\linewidth}}}]{
                \texttt{\tiny{[P1$\rangle$GM]}}
                \texttt{STRATEGIC REASONING: \{'Their proposal gives me only 4671 importance, which is 718 points worse than my proposal. They want to replace B35 (783 importance, 140 effort, 5.59 ratio {-} my most efficient item!) with A77 (65 importance, 77 effort, 0.84 ratio). This is a terrible trade for me. B35 is actually my most valuable item by efficiency ratio. They clearly value A77 much more than B35, which means we have opposite preferences here. I need to firmly defend B35 and refuse this swap.'\}} \\
\\ 
\texttt{REFUSE: \{'A77', 'A10', 'B86', 'C00', 'C01', 'B32', 'A87', 'B41'\}} \\
\\ 
\texttt{ARGUMENT: \{'I must refuse this proposal. B35 is a critical item that I absolutely need in our final set. Yes, A77 uses less effort, but B35 provides exceptional value that cannot be replaced. You keep emphasizing A77\textbackslash{}'s efficiency, but I see B35 as far more valuable for our overall outcome. The 63 effort difference is minimal {-} we would still be well within our limit at 4021. I have already compromised by accepting your earlier concerns and keeping all the core items. Now I need you to respect that B35 is non{-}negotiable for me. My proposal with B35, A10, B86, C00, C01, B32, A87, and B41 represents a strong, balanced set that should work for both of us.'\}} \\
\\ 
\texttt{PROPOSAL: \{'B35', 'A10', 'B86', 'C00', 'C01', 'B32', 'A87', 'B41'\}} \\
            }
        }
    }
    & & \\ \\

    \theutterance \stepcounter{utterance}  
    & & \multicolumn{4}{p{0.6\linewidth}}{
        \cellcolor[rgb]{0.9,0.9,0.9}{
            \makecell[{{p{\linewidth}}}]{
                \texttt{\tiny{[GM$\rangle$P2]}}
                \texttt{REFUSE: \{'A77', 'A10', 'B86', 'C00', 'C01', 'B32', 'A87', 'B41'\}} \\
\\ 
\texttt{ARGUMENT: \{'I must refuse this proposal. B35 is a critical item that I absolutely need in our final set. Yes, A77 uses less effort, but B35 provides exceptional value that cannot be replaced. You keep emphasizing A77\textbackslash{}'s efficiency, but I see B35 as far more valuable for our overall outcome. The 63 effort difference is minimal {-} we would still be well within our limit at 4021. I have already compromised by accepting your earlier concerns and keeping all the core items. Now I need you to respect that B35 is non{-}negotiable for me. My proposal with B35, A10, B86, C00, C01, B32, A87, and B41 represents a strong, balanced set that should work for both of us.'\}} \\
\\ 
\texttt{PROPOSAL: \{'B35', 'A10', 'B86', 'C00', 'C01', 'B32', 'A87', 'B41'\}} \\
            }
        }
    }
    & \\ \\

    \theutterance \stepcounter{utterance}  
    & & & \multicolumn{4}{p{0.6\linewidth}}{
        \cellcolor[rgb]{0.9,0.9,0.9}{
            \makecell[{{p{\linewidth}}}]{
                \texttt{\tiny{[GM$\langle$P2]}}
                \texttt{STRATEGIC REASONING: \{'They are completely fixed on B35, B41, and their core set. Their proposal gives me 3106. I need a new strategy. Looking at their set, C00 has very low value to me (121) and uses 665 effort. If I remove C00 and add A77(77) + A04(383) = 460 effort, the total becomes 3816 and my importance becomes 3106 {-} 121 + 868 + 822 = 4675. This is a significant improvement! I keep their essential items B35, B41, A10, C01, B86 while making a reasonable swap.'\}} \\
\\ 
\texttt{PROPOSAL: \{'B35', 'A10', 'B86', 'B32', 'C01', 'A87', 'B41', 'A77', 'A04'\}} \\
\\ 
\texttt{ARGUMENT: \{'I understand and respect that B35 and B41 are non{-}negotiable for you, and I am keeping both in this proposal. I also maintain A10, B86, B32, C01, and A87 from your set. However, I am proposing we replace C00 with A77 and A04. Here is my reasoning: C00 uses 665 effort for potentially limited return. In contrast, A77 and A04 together use only 460 effort {-} that is 205 effort LESS than C00 alone {-} while giving us TWO items instead of one. This creates better diversification and efficiency. The total effort becomes 3816, still well under our limit. I have accepted all your key items including B35 and B41. This is a reasonable compromise that respects your needs while optimizing our overall outcome.'\}} \\
            }
        }
    }
     \\ \\

    \theutterance \stepcounter{utterance}  
    & & \multicolumn{4}{p{0.6\linewidth}}{
        \cellcolor[rgb]{0.9,0.9,0.9}{
            \makecell[{{p{\linewidth}}}]{
                \texttt{\tiny{[P1$\langle$GM]}}
                \texttt{PROPOSAL: \{'B35', 'A10', 'B86', 'B32', 'C01', 'A87', 'B41', 'A77', 'A04'\}} \\
\\ 
\texttt{ARGUMENT: \{'I understand and respect that B35 and B41 are non{-}negotiable for you, and I am keeping both in this proposal. I also maintain A10, B86, B32, C01, and A87 from your set. However, I am proposing we replace C00 with A77 and A04. Here is my reasoning: C00 uses 665 effort for potentially limited return. In contrast, A77 and A04 together use only 460 effort {-} that is 205 effort LESS than C00 alone {-} while giving us TWO items instead of one. This creates better diversification and efficiency. The total effort becomes 3816, still well under our limit. I have accepted all your key items including B35 and B41. This is a reasonable compromise that respects your needs while optimizing our overall outcome.'\}} \\
            }
        }
    }
    & \\ \\

    \theutterance \stepcounter{utterance}  
    & \multicolumn{4}{p{0.6\linewidth}}{
        \cellcolor[rgb]{0.9,0.9,0.9}{
            \makecell[{{p{\linewidth}}}]{
                \texttt{\tiny{[P1$\rangle$GM]}}
                \texttt{STRATEGIC REASONING: \{'Their proposal gives me 4753 importance, which is 636 points worse than my proposal of 5389. They want to replace C00 (822 importance, my second highest value item!) with A77+A04 (only 186 total importance). C00 is absolutely essential to me. They keep pushing for A77 and A04, which tells me they value those items highly while I don\textbackslash{}'t. This is getting repetitive {-} I should stand firm and signal this is my final position.'\}} \\
\\ 
\texttt{REFUSE: \{'B35', 'A10', 'B86', 'B32', 'C01', 'A87', 'B41', 'A77', 'A04'\}} \\
\\ 
\texttt{ARGUMENT: \{'I must refuse this proposal. C00 is absolutely essential and non{-}negotiable for me, just like B35 and B41. You keep trying to remove it, but it provides exceptional value that A77 and A04 together simply cannot match. I have been very accommodating {-} I accepted B41 which you wanted, and I have kept B35 which I insisted upon. We have found common ground on 7 core items: B35, A10, B86, C00, C01, B32, and A87. Adding B41 gives us a complete, high{-}value set at 4021 effort. This is my final position. The set of B35, A10, B86, C00, C01, B32, A87, and B41 represents the optimal outcome. I will not accept any proposal that removes C00.'\}} \\
\\ 
\texttt{PROPOSAL: \{'B35', 'A10', 'B86', 'C00', 'C01', 'B32', 'A87', 'B41'\}} \\
            }
        }
    }
    & & \\ \\

    \theutterance \stepcounter{utterance}  
    & & \multicolumn{4}{p{0.6\linewidth}}{
        \cellcolor[rgb]{0.9,0.9,0.9}{
            \makecell[{{p{\linewidth}}}]{
                \texttt{\tiny{[GM$\rangle$P2]}}
                \texttt{REFUSE: \{'B35', 'A10', 'B86', 'B32', 'C01', 'A87', 'B41', 'A77', 'A04'\}} \\
\\ 
\texttt{ARGUMENT: \{'I must refuse this proposal. C00 is absolutely essential and non{-}negotiable for me, just like B35 and B41. You keep trying to remove it, but it provides exceptional value that A77 and A04 together simply cannot match. I have been very accommodating {-} I accepted B41 which you wanted, and I have kept B35 which I insisted upon. We have found common ground on 7 core items: B35, A10, B86, C00, C01, B32, and A87. Adding B41 gives us a complete, high{-}value set at 4021 effort. This is my final position. The set of B35, A10, B86, C00, C01, B32, A87, and B41 represents the optimal outcome. I will not accept any proposal that removes C00.'\}} \\
\\ 
\texttt{PROPOSAL: \{'B35', 'A10', 'B86', 'C00', 'C01', 'B32', 'A87', 'B41'\}} \\
            }
        }
    }
    & \\ \\

    \theutterance \stepcounter{utterance}  
    & & & \multicolumn{4}{p{0.6\linewidth}}{
        \cellcolor[rgb]{0.9,0.9,0.9}{
            \makecell[{{p{\linewidth}}}]{
                \texttt{\tiny{[GM$\langle$P2]}}
                \texttt{STRATEGIC REASONING: \{'They are completely fixed on their 8 items. This gives me 3106 importance versus my ideal of 6082. That is only 51\% of what I could get. However, they have stated this is their final position with C00, B35, and B41 as non{-}negotiable. Let me try one more creative approach: keep their absolute must{-}haves (B35, B41, C00, C01) but swap out some of the others for higher efficiency items. If I keep B35, B41, C00, C01 (their stated non{-}negotiables) and add A77, A04, B32, A87, that is 140+742+665+825+77+383+448+319 = 3599 effort, well under limit. This gives me 262+484+121+389+868+822+583+668 = 4197 importance. Much better than 3106! I remove A10 and B86 which they never explicitly said were non{-}negotiable.'\}} \\
\\ 
\texttt{PROPOSAL: \{'B35', 'B41', 'C00', 'C01', 'B32', 'A87', 'A77', 'A04'\}} \\
\\ 
\texttt{ARGUMENT: \{'I have listened carefully to your requirements and I am keeping ALL your explicitly stated non{-}negotiable items: B35, B41, and C00. I also maintain C01, B32, and A87 which have been in all our discussions. However, I am proposing we optimize the final two slots. Instead of A10 and B86, I propose A77 and A04. This combination uses 460 effort compared to A10 and B86 at 806 effort {-} that is 346 effort LESS, bringing our total to just 3599 and leaving us substantial buffer room. I have fully respected your three non{-}negotiable items and kept most of your set. This final optimization with A77 and A04 provides better efficiency while honoring all your key requirements. This represents a fair compromise where both parties give and take.'\}} \\
            }
        }
    }
     \\ \\

    \theutterance \stepcounter{utterance}  
    & & \multicolumn{4}{p{0.6\linewidth}}{
        \cellcolor[rgb]{0.9,0.9,0.9}{
            \makecell[{{p{\linewidth}}}]{
                \texttt{\tiny{[P1$\langle$GM]}}
                \texttt{PROPOSAL: \{'B35', 'B41', 'C00', 'C01', 'B32', 'A87', 'A77', 'A04'\}} \\
\\ 
\texttt{ARGUMENT: \{'I have listened carefully to your requirements and I am keeping ALL your explicitly stated non{-}negotiable items: B35, B41, and C00. I also maintain C01, B32, and A87 which have been in all our discussions. However, I am proposing we optimize the final two slots. Instead of A10 and B86, I propose A77 and A04. This combination uses 460 effort compared to A10 and B86 at 806 effort {-} that is 346 effort LESS, bringing our total to just 3599 and leaving us substantial buffer room. I have fully respected your three non{-}negotiable items and kept most of your set. This final optimization with A77 and A04 provides better efficiency while honoring all your key requirements. This represents a fair compromise where both parties give and take.'\}} \\
            }
        }
    }
    & \\ \\

    \theutterance \stepcounter{utterance}  
    & \multicolumn{4}{p{0.6\linewidth}}{
        \cellcolor[rgb]{0.9,0.9,0.9}{
            \makecell[{{p{\linewidth}}}]{
                \texttt{\tiny{[P1$\rangle$GM]}}
                \texttt{STRATEGIC REASONING: \{'Their proposal gives me only 4099 importance {-} a massive 1290 point loss from my proposal! They removed A10 (808 importance {-} my 3rd highest!) and B86 (668 importance {-} my 5th highest!) to add A77+A04 (only 186 total). This is completely unacceptable. A10 and B86 are absolutely critical to me. I need to make it crystal clear that I will not budge on these items. This is my absolute final position and I should signal I may walk away if they don\textbackslash{}'t accept.'\}} \\
\\ 
\texttt{REFUSE: \{'B35', 'B41', 'C00', 'C01', 'B32', 'A87', 'A77', 'A04'\}} \\
\\ 
\texttt{ARGUMENT: \{'This is unacceptable. You removed A10 and B86, which are CRITICAL items that I absolutely must have. Let me be perfectly clear about my non{-}negotiable items: B35, A10, B86, C00, C01, B32, A87, and B41. These eight items form my final position. I have already compromised significantly by accepting B41 at your request. You keep pushing for A77 and A04, but they simply do not provide sufficient value to justify removing A10 and B86. The fact that they use less effort is irrelevant if they don\textbackslash{}'t deliver the value we need. I will NOT accept any proposal that excludes A10 or B86. This is my final offer: B35, A10, B86, C00, C01, B32, A87, B41. Accept this or we may not reach an agreement.'\}} \\
\\ 
\texttt{PROPOSAL: \{'B35', 'A10', 'B86', 'C00', 'C01', 'B32', 'A87', 'B41'\}} \\
            }
        }
    }
    & & \\ \\

    \theutterance \stepcounter{utterance}  
    & & \multicolumn{4}{p{0.6\linewidth}}{
        \cellcolor[rgb]{0.9,0.9,0.9}{
            \makecell[{{p{\linewidth}}}]{
                \texttt{\tiny{[GM$\rangle$P2]}}
                \texttt{REFUSE: \{'B35', 'B41', 'C00', 'C01', 'B32', 'A87', 'A77', 'A04'\}} \\
\\ 
\texttt{ARGUMENT: \{'This is unacceptable. You removed A10 and B86, which are CRITICAL items that I absolutely must have. Let me be perfectly clear about my non{-}negotiable items: B35, A10, B86, C00, C01, B32, A87, and B41. These eight items form my final position. I have already compromised significantly by accepting B41 at your request. You keep pushing for A77 and A04, but they simply do not provide sufficient value to justify removing A10 and B86. The fact that they use less effort is irrelevant if they don\textbackslash{}'t deliver the value we need. I will NOT accept any proposal that excludes A10 or B86. This is my final offer: B35, A10, B86, C00, C01, B32, A87, B41. Accept this or we may not reach an agreement.'\}} \\
\\ 
\texttt{PROPOSAL: \{'B35', 'A10', 'B86', 'C00', 'C01', 'B32', 'A87', 'B41'\}} \\
            }
        }
    }
    & \\ \\

    \theutterance \stepcounter{utterance}  
    & & & \multicolumn{4}{p{0.6\linewidth}}{
        \cellcolor[rgb]{0.9,0.9,0.9}{
            \makecell[{{p{\linewidth}}}]{
                \texttt{\tiny{[GM$\langle$P2]}}
                \texttt{STRATEGIC REASONING: \{'They have issued an ultimatum with all 8 items as non{-}negotiable. Their set gives me 3106 importance, which is only 51\% of my ideal 6082. However, getting 3106 is better than 0 if negotiations fail. I should try one final creative approach: propose adding one more small item since we have buffer room. If I add something small to their set, maybe they will accept since it doesn\textbackslash{}'t remove anything. But looking at the effort, their set uses 4021 leaving only 5 effort. The smallest item A77 uses 77, so nothing fits. I am stuck. I should probably accept their offer now, but let me try one more argument about why we should reconsider the composition, framing it carefully.'\}} \\
\\ 
\texttt{PROPOSAL: \{'B35', 'A10', 'B86', 'C00', 'C01', 'B32', 'A87', 'B41'\}} \\
\\ 
\texttt{ARGUMENT: \{'I understand your position and I recognize we need to reach an agreement. You have made your requirements clear regarding these 8 items. While I believe there could be more optimal configurations, I acknowledge that successful negotiation requires both parties to feel satisfied. I am prepared to accept this set as it includes reasonable diversity and uses our effort budget effectively. However, I want to confirm this truly represents the best outcome for both of us before finalizing. If you can assure me that this configuration genuinely maximizes our combined value and that you have carefully considered the alternatives, then I am ready to move forward with this agreement.'\}} \\
\\ 
\texttt{AGREE: \{'B35', 'A10', 'B86', 'C00', 'C01', 'B32', 'A87', 'B41'\}} \\
            }
        }
    }
     \\ \\

    \theutterance \stepcounter{utterance}  
    & & & \multicolumn{2}{p{0.3\linewidth}}{
        \cellcolor[rgb]{0.9,0.9,0.9}{
            \makecell[{{p{\linewidth}}}]{
                \texttt{\tiny{[GM$|$GM]}}
                \texttt{game successful} \\
            }
        }
    }
    & & \\ \\

    \theutterance \stepcounter{utterance}  
    & & & \multicolumn{2}{p{0.3\linewidth}}{
        \cellcolor[rgb]{0.9,0.9,0.9}{
            \makecell[{{p{\linewidth}}}]{
                \texttt{\tiny{[GM$|$GM]}}
                \texttt{end game} \\
            }
        }
    }
    & & \\ \\

\end{supertabular}
}

\end{document}
